\scanpage{56}
% archive sheet number: 55
Algèbres étales et la différente \mspage{(2)}\hfill Tribu n\textsuperscript{o} 65

\uline{n\textsuperscript{o} 2}. \uline{Algèbres étales et algèbres non ramifiées}.

Soit $B$ une $A$-algèbre. On considère les propriétés suivantes :

a)\quad $B$ est de p. f.

b)\quad $B$ est plate

c) Pour tout $\handwritten{\unsure{q|\mathfrak{q}}\in}\operatorname{Spec} B$ et pour $\handwritten{\unsure{p|\mathfrak{p}}}= \handwritten{q\cap}A$, $B_{\handwritten{q}}/\handwritten{p}B_{\handwritten{q}}$ est une extension
séparable finie de $A_{\handwritten{p}}/\handwritten{p}A_{\handwritten{p}}$.

En présence de a) \deleted{et b)}, c) veut dire que $\handwritten{\Omega}_{B/A}$ est nul en $\handwritten{q}$.

On dit que $B$ est étale si on a a), b), c).

On dit que $B$ est étale en $\handwritten{q\in}\operatorname{Spec}(B)$ s'il existe $f\handwritten{\notin q}$ tel que on ait a),
b), c) pour $B_f$. (i. e. tel que $B_f$ soit étale).

On dit que $B$ est non ramifiée si elle est de p. f. et si $\handwritten{\Omega}_{B/A}=0$.

Ainsi : étale $=$ non ramifiée $+$ plate.

Si $A$ est réduite, $B$ est réduite.

Transitivité ($A\longrightarrow B\longrightarrow C$) de l'étalité et de la non-ramification.

Effet d'un changement d'anneau sur ces notions : la montée marche dans le cas
général, la descente marche dans le cas fidèlement plat.

\uline{n\textsuperscript{o} 3}. \uline{Algèbres finies étales et non ramifiées}.

Les équivalences de la prop. 1, p. 1, du texte (améliorées comme dans la Tribu
n\textsuperscript{o} 64 ; on supposera seulement $B$ de présentation finie sur $A$ ).

Soient $B$ de p. f. sur $A$ noethérien, $\handwritten{q\in}\operatorname{Spec} B$ et $\handwritten{p}=\handwritten{q\cap} A$ ; pour
que $B$ soit étale en $\handwritten{q}$, il faut et il suffit que $\handwritten{\widehat}{B}_{\handwritten{q}}$ soit étale finie sur
$\handwritten{\widehat}{A}_{\handwritten{p}}$ ; analogue non ramifié ; la démonstration utilise la platitude supérieure. Corollaire pour le cas où les corps résiduels sont égaux (cf. S. G. A.).

Pour que $B$ soit étale finie sur $A$ connexe, il faut et il suffit qu'il existe
$A'$ fidèlement plate sur $A$ telle que $A'\handwritten{\otimes}B$ soit un produit de copies de $A'$
(si $A$ n'est pas connexe, ça ne marche que par morceaux).

La prop. 5 du texte (arrangée comme dans la Tribu n\textsuperscript{o} 64). Le calcul de traces de
la prop. 3, a) du \S~2 (p. 16).

\uline{n\textsuperscript{o} 4}. \uline{Structure locale des algèbres non ramifiées et étales}.

Soit $B$ une $A$-algèbre finie et soit $\handwritten{q\in}\operatorname{Spec}(B)$. Pour que $B$ soit non
ramifiée en $\handwritten{q}$ il faut et il suffit qu'il existe $x\handwritten{\in} B'$ de polynôme minimal $F$
